\subsection{Counting Subtrees}

\begin{claim}
  \label{claim:counting_trees}
  Let $G=(V,E)$ be a graph of maximum degree at most $(\Delta+1)$, and let $k$ be an
  integer less than $|E|$. Fix a vertex $r\in V$. There is a constant
  $\alpha$, depending only on $\Delta$, such that the number of
  \emph{connected} subgraphs containing $r$ and having fewer than $k$
  edges is at most $\alpha^k$.
\end{claim}

\begin{proof}
  If $k\le 0$, there are no such subgraphs, so the claim is immediate.
  Set $d=\Delta+1$. For each connected subgraph $H$ containing $r$ and
  having exactly $j$ edges, duplicate every edge of $H$. The resulting
  connected multigraph has even degrees, and hence admits an Eulerian
  circuit starting at $r$. This circuit induces a closed walk of length
  $2j$ in $G$ that traverses precisely the edges of $H$ (each twice).

  Fix an ordering of the edges of $G$ and choose, for each $H$, one such
  walk deterministically, for example the lexicographically first one.
  The correspondence is injective: from the selected walk we recover
  $H$ as the subgraph consisting of all edges traversed by the walk,
  together with $r$ in the case $j=0$.

  Since every vertex has degree at most $d$, there are at most $d^{2j}$ walks
  of length $2j$ starting at $r$. Consequently, the required number
  of subgraphs is at most
  \[
    \sum_{j=0}^{k-1}d^{2j}
    \le \sum_{j=0}^{k-1}(d+1)^{2j}
    \le (d+1)^{2k}.
  \]
  Thus one may take $\alpha=(\Delta+2)^2$, which depends only on
  $\Delta$.
\end{proof}
